\section{What old common-graph primitives do imply for one reduced repair}
\label{sec:136-old-graph-linking-repair}

\paragraph{Archive and exact scope.}
This concerns cote136, physical pp.~117,119,121,
footers116,118,120, scans117,119,121.  It gives a positive
one-stage application of ARCH-136-LINKING-CHAIN-GATE under
a stronger hypothesis than old degree-two vertex compatibility.
The hypothesis is the old degree-three common-graph primitive
property itself.  The result does not prove that this stronger
property is preserved by the new stage.  The complementary
LOW-FACTOR-PRESERVATION theorem gives one sufficient invariant
for particular further operations; its degree-four stop remains.

Put \(S=\Gamma_{\mathbb Z}(T)\),
\(J=\mathbb Q[T]/S\).  Let \(A\rightrightarrows C_j\)
be finitely many strict maps to admissible targets with literal
common first skeleton \(F\); all maps are the identity on
\(F\), and \(F_0=S\).  The algebra \(F_1\) is the
specified free even zero-form unit-disk algebra, so
\(H^0(F_1)=S\), \(H^i(F_1)=0\) for \(i>0\), in
every weight.  Assume in every simplex and weight that
\[
 H^0(A)=S,\qquad H^1(A)=0,\qquad
                         H^2(A)\text{ is integer torsion}.
                                                               \tag{136.OG1}
\]
Choose any set of old torsion pairs
\[
 z_\alpha\in A^{v_\alpha,3}(\Delta[n_\alpha]),\quad
 e_\alpha\in A^{v_\alpha,2}(\Delta[n_\alpha]),\quad
 dz_\alpha=0,\quad de_\alpha=r_\alpha z_\alpha,
              \qquad n_\alpha\geq1,\quad r_\alpha\geq1.
                                                               \tag{136.OG2}
\]
Their vertex values are zero because \(A_0=S\).

The stronger old hypothesis is that for each of these pairs
there are target sections
\[
 dC_{j\alpha}=f_j(z_\alpha),\qquad
 dT_{j\alpha}=f_j(e_\alpha)-r_\alpha C_{j\alpha},
       \qquad |C_{j\alpha}|=2,\quad |T_{j\alpha}|=1,
                                                               \tag{136.OG3}
\]
whose restrictions to
\(G_\alpha=\operatorname{sk}_1\Delta[n_\alpha]\)
are identical for every target, as actual sections of \(F(G_\alpha)\).
It suffices to assume the all-modulus common-graph primitive
property for \(A\) in a fixed tail, and take the weights
\(v_\alpha\) in that tail.  In addition take them above
one fixed target bound for vertical comparison through degree
two and surjectivity of the form-zero restriction to all vertices.
Both bounds are independent of \(n_\alpha,r_\alpha\), and
the number of pairs.

Adjoin the actual reduced-representable pairs
\[
 A\longrightarrow E,
 \qquad |c_\alpha|=2,\quad |t_\alpha|=1,
 \qquad dc_\alpha=z_\alpha,\quad
              dt_\alpha=e_\alpha-r_\alpha c_\alpha,
                                                               \tag{136.OG4}
\]
with new generators on
\(\Delta[n_\alpha]/\operatorname{sk}_0\Delta[n_\alpha]\)
and zero vertex values.  Suppose arbitrary target extensions
\(g_j:E\to C_j\) have already been chosen.

\paragraph{ARCH-136-OLD-GRAPH-LINKING-REPAIR.}
Under these hypotheses the linking difference of the arbitrary
extensions satisfies the exact chain condition (136.LK16).
It can be made zero, for every finite modulus in one uniform
output tail, by changing only the images of the new \(t_\alpha\)
by closed one-forms.  All old maps and all new \(c_\alpha\)
images are retained.  Alternatively one can use (136.OG3) to
attach the pairs relative to the entire graph; this retains
the literal first skeleton \(F\) and all degree-two vertex
compatibilities.  Neither assertion includes preservation of
the old degree-three common-graph property after that stage.

\begin{proof}
The mixed-term calculation (136.LK8)--(136.LK9) applies to
these actual free pairs as well.  It yields
\[
 H^0(E)=S,\quad H^1(E)=0,\qquad
 0\longrightarrow H^2(A)\longrightarrow H^2(E)
                         \longrightarrow K\longrightarrow0,
                                                               \tag{136.OG5}
\]
where
\[
 \begin{split}
 L&=\bigoplus_\alpha(S/r_\alpha)\langle v_\alpha\rangle
       \otimes_{\mathbb Z}
       \mathbb Z[\Delta[n_\alpha]]/
                         \mathbb Z[\operatorname{sk}_0\Delta[n_\alpha]],\\
 K&=\ker\{L\longrightarrow H^3(A),\ c_\alpha\longmapsto[z_\alpha]\}.
 \end{split}                                                \tag{136.OG6}
\]
All the assertions are simplicially natural.  In particular
\(H^2(E)\) is integer torsion, so the universal linking
differences (136.LK10) apply in the fixed target comparison
range, simultaneously for every modulus.

First construct reference extensions \(g_j^\flat\) by
assigning \(c_\alpha,t_\alpha\) the images
\(C_{j\alpha},T_{j\alpha}\) from (136.OG3).  The free
reduced presentation gives genuine strict simplicial maps.
They agree on the entire new first skeleton \(E_1\),
although \(E_1\) need not itself equal \(F_1\).

Their vertex linking differences are zero.  To check this
without presuming exactness of \(E_1\), take a torsion pair
\(dx=0,\ db=Nx\) of forms two and one in \(E_n\).
On each edge its images under all \(g_j^\flat\) are the
same pair in \(F_1\).  Choose there a one-form primitive
\(a_e\) of the two-form and a zeroform primitive \(q_e\)
of the one-form minus \(Na_e\).  These exist by the
all-weight vertical exactness of \(F_1\).
Any target choices \(a_j,q_j\) used in (136.LK10) differ
on that edge by
\[
             a_j|_e-a_e=dh_{je},\qquad
             q_j|_e-q_e+Nh_{je}\text{ is a coefficient}.
                                                               \tag{136.OG7}
\]
Thus their endpoint differences agree modulo \(N\) on
every edge.  The complete graph on the simplex vertices is
connected, so their whole vertex vectors agree modulo the
diagonal and \(N\).  This proves vanishing of the universal
linking difference for \(g^\flat\), in every weight where
the target choices are defined.  The same argument applies
to the old maps on \(A\).

Compare an arbitrary extension \(g_j\) with \(g_j^\flat\).
On a generating cell put
\[
 \delta c_{j\alpha}=g_j(c_\alpha)-C_{j\alpha},\qquad
 \delta t_{j\alpha}=g_j(t_\alpha)-T_{j\alpha}.
                                                               \tag{136.OG8}
\]
They satisfy \(d\delta c=0\) and
\(d\delta t=-r_\alpha\delta c\).  The common fixed
weight bound supplies
\[
 da_{j\alpha}=\delta c_{j\alpha},\qquad
 dq_{j\alpha}=\delta t_{j\alpha}+r_\alpha a_{j\alpha},
             \qquad |a_{j\alpha}|=1,\quad |q_{j\alpha}|=0.
                                                               \tag{136.OG9}
\]
Their vertex vector defines a canonical element
\[
 \lambda_{j\alpha}
       =\left[\frac{q_{j\alpha}|_{\mathrm{vertices}}}{r_\alpha}\right]
       \in\mathcal V_{n_\alpha}(J)^{v_\alpha}[r_\alpha],
             \qquad \mathcal V_n(J)=J^{n+1}/\operatorname{diag}J.
                                                               \tag{136.OG10}
\]
Changing \(a\) by an exact one-form changes \(q/r\)
by an integral zeroform at the vertices; a change of its
closed zeroform ambiguity is diagonal.  Target \(H^1=0\)
and \(H^0=S\) therefore show independence of all choices.
The displayed element is killed by \(r_\alpha\), since
the original vertex vector is integral.  The universal property
of the reduced representables extends these elements to a
weight-preserving simplicial coefficient map
\[
                            \lambda_j:L\longrightarrow\mathcal V(J).
                                                               \tag{136.OG11}
\]
This uses the actual cells of all dimensions; it does not
assume that \(NL\) has zero differential.

We verify its relation to the linking difference, with all
orders and signs explicit.  Represent a two-class as
\(x=a+\sum_i s_i c_i\) in one simplex, with \(s_i\in S\),
as in (136.LK9).  Choose a common multiple \(N\) of the
relevant \(r_i\) and the class annihilator, enlarging it
to kill a remaining old two-class if necessary.  A relation
can then be chosen in the form
\[
                db=Nx,
                \qquad b=b_{\mathrm{old}}
                              -\sum_i(N/r_i)s_i t_i.
                                                               \tag{136.OG12}
\]
The old torsion hypothesis in (136.OG1) is exactly what makes
that last enlargement possible.  Between \(g_j\) and
\(g_j^\flat\), the target image of \(x\) changes by
\(d(\sum_i s_i a_{ji})\), while that of \(b\) changes
by \(-\sum_i(N/r_i)s_i\delta t_{ji}\).  Choosing its
one-form primitive accordingly, the closed one-form
\(g_j(b)-N a_j\) changes by
\[
                   -d\sum_i(N/r_i)s_i q_{ji}.
                                                               \tag{136.OG13}
\]
After division by \(N\) this changes the linking value by
\(-\sum_i s_i\lambda_{ji}\).  Taking target \(0\)
as reference and using the zero linking difference of
\(g^\flat\), the arbitrary extensions therefore satisfy
\[
             \Delta_j^{g}
                        =(-\lambda_j+\lambda_0)|_K
                 \quad\text{in the specified output tail}.
                                                               \tag{136.OG14}
\]
All computations commute with simplicial pullbacks.  Formula
(136.OG14) is the desired explicit simplicial extension to
\(L\).  In particular its normalized degree-one component
vanishes on \(N_1K\cap\partial(N_2L)\), proving the
exact condition (136.LK16).

It also supplies a concrete correction retaining every new
two-form image.  On each generating cell choose a target
zeroform \(h_{j\alpha}\) whose vertex vector is
\[
        -q_{j\alpha}|_{\mathrm{vertices}}
                              +q_{0\alpha}|_{\mathrm{vertices}},
                                                               \tag{136.OG15}
\]
identifying these coefficient values in the common \(S\).
The fixed form-zero row bound gives the extension from all
vertices in every simplicial dimension.  Set
\(g_j(t_\alpha)\gets g_j(t_\alpha)+d h_{j\alpha}\),
and leave every \(c_\alpha\) and old image unchanged.
This preserves the differential equations and zero vertex
values, so the free reduced pushout gives strict maps under
\(A\).  The computation (136.OG12)--(136.OG13), applied
to this change, subtracts exactly
\(-\lambda_j+\lambda_0\) from the linking difference.
Consequently all finite-modulus vertex residues agree after
the correction, with the same one output cutoff.  No upper bound
on the orders, number of cells, or their weights has been used.

Finally, the common restrictions in (136.OG3) are actual
elements of the old graph algebra \(A(G_\alpha)=F(G_\alpha)\).
One may instead identify the whole graph of each new pair
with these old sections and attach relative to
\(\operatorname{sk}_1\Delta[n_\alpha]\).
All differential equations hold there, so this is a genuine
strict PD pushout with the reference target extensions.
Its first skeleton is literally \(F\).  The same edge
argument (136.OG7) gives all degree-two vertex compatibilities.
This alternative changes the source attachment, rather than
claiming to retain arbitrary initially prescribed target images.
\end{proof}

\paragraph{Exact consequence and remaining gate.}
Old degree-three common-graph primitives force the linking
chain condition for the subsequent reduced degree-three stage;
old degree-two vertex compatibility by itself does not, as
LINKING-TRIANGLE-OBSTRUCTION shows.  The conclusion is uniform
in every finite order and uses one bound on the actual cell
weights.  A full iteration still needs a reason that the
stronger old common-graph condition survives each required
vertical repair.  This module supplies no such unrestricted
iteration and leaves the integral common-source problem open.

\paragraph{Verification status.}
The author completed a whole saved-proof reread: actual reduced
PD presentation and low cohomology; reference maps and edge
normalization; independence and simplicial extension of every
coefficient linking value; the comparison and correction signs;
all finite moduli with one cutoff; and the graph-relative
alternative.  Root and reviewer15 have completed whole saved-proof checks
with PASS.  Reviewer15 contributed an alternative triangle
calculation; no independent rediscovery is claimed.  The
stronger graph property is not claimed to survive this stage.
